\documentclass[12pt,a4paper]{article}
\usepackage[margin=1in]{geometry}
\usepackage{graphicx}
\usepackage{listings}
\usepackage{xcolor}
\usepackage{booktabs}
\usepackage{hyperref}
\usepackage{fancyhdr}
\usepackage{titlesec}
\usepackage{caption}

\definecolor{codegreen}{rgb}{0,0.5,0}
\definecolor{codegray}{rgb}{0.5,0.5,0.5}
\definecolor{codebg}{rgb}{0.97,0.97,0.97}
\definecolor{codeblue}{rgb}{0.1,0.1,0.7}

\lstdefinestyle{verilog}{
    backgroundcolor=\color{codebg},
    commentstyle=\color{codegreen},
    keywordstyle=\color{codeblue}\bfseries,
    numberstyle=\tiny\color{codegray},
    basicstyle=\ttfamily\footnotesize,
    breaklines=true,
    frame=single,
    numbers=left,
    numbersep=5pt,
    tabsize=4,
    morekeywords={module,endmodule,input,output,wire,reg,assign,always,begin,end,case,endcase,if,else,posedge,default,parameter,initial}
}

\lstdefinestyle{asm}{
    backgroundcolor=\color{codebg},
    commentstyle=\color{codegreen},
    basicstyle=\ttfamily\footnotesize,
    breaklines=true,
    frame=single,
    numbers=left,
    numbersep=5pt,
    tabsize=4,
}

\pagestyle{fancy}
\fancyhf{}
\rhead{Computer Architecture Project}
\lhead{EE/CE 321L/371L}
\rfoot{Page \thepage}

\title{
    \vspace{-1cm}
    \textbf{Course Project Report} \\
    \large Assembly-Level Execution on Integrated RISC-V Processor \\
    \vspace{0.3cm}
    \normalsize EE/CE 321L/371L -- Computer Architecture Lab
}
\author{
    Name: \underline{\hspace{4cm}} \quad
    ID: \underline{\hspace{3cm}} \quad
    Section: \underline{\hspace{2cm}}
}
\date{\today}

\begin{document}
\maketitle
\thispagestyle{fancy}

% ============================================================
\section{Part A: Base Assembly Program Execution}
% ============================================================

\subsection{Objective}
Execute the RISC-V assembly program developed in Lab 10 (FSM-based countdown using switches and LEDs) on the integrated single-cycle RISC-V processor from Lab 11.

\subsection{Program Description}
The Lab 10 program implements a Finite State Machine with two states:

\begin{itemize}
    \item \textbf{WAIT state:} The processor continuously reads the switch value from the memory-mapped switch address (768 = \texttt{0x300}). If the value is zero, it loops. When a non-zero value is detected, it transitions to the COUNT state.
    \item \textbf{COUNT state:} The captured value is passed to a countdown subroutine. The subroutine displays the current count on the LEDs (address 512 = \texttt{0x200}), decrements by 1, runs a small delay loop, and repeats until the count reaches zero. It then clears the LEDs and returns to the WAIT state.
\end{itemize}

The subroutine uses proper stack conventions --- it saves the return address (\texttt{ra}) and the argument register (\texttt{a2}) onto the stack before execution, and restores them before returning via \texttt{jalr}.

\subsection{Assembly Listing}
\begin{lstlisting}[style=asm, caption={Lab 10 FSM Program}]
# Initialization
addi x28, x0, 5       # test value
addi x2, x0, 511      # sp = 511
addi x5, x0, 768      # switch address
addi x6, x0, 512      # LED address
sw   x28, 0(x5)        # write test value to switches
sw   x0, 0(x6)         # clear LEDs

# WAIT LOOP
wait: lw x11, 0(x5)    # read switch value
      beq x11, x0, wait # if zero, keep waiting

# Call countdown subroutine
add  x10, x11, x0      # a0 = switch value
jal  x1, countdown      # call subroutine
beq  x0, x0, clear      # after return, jump to clear

# COUNTDOWN SUBROUTINE
countdown:
    addi x2, x2, -8     # push stack
    sw   x1, 4(x2)      # save ra
    sw   x12, 0(x2)     # save a2
    add  x12, x10, x0   # a2 = argument
loop:
    sw   x12, 0(x6)     # display on LEDs
    beq  x12, x0, done  # if zero, done
    addi x12, x12, -1   # decrement
    addi x13, x0, 3     # delay counter
delay:
    addi x13, x13, -1   # delay--
    bne  x13, x0, delay # delay loop
    beq  x0, x0, loop   # back to display
done:
    sw   x0, 0(x6)      # clear LEDs
    lw   x12, 0(x2)     # restore a2
    lw   x1, 4(x2)      # restore ra
    addi x2, x2, 8      # pop stack
    jalr x0, 0(x1)      # return

# HALT
halt: jal x0, halt      # infinite loop
\end{lstlisting}

\subsection{Instructions Used}
\texttt{ADDI}, \texttt{SW}, \texttt{LW}, \texttt{ADD}, \texttt{BEQ}, \texttt{BNE}, \texttt{JAL}, \texttt{JALR}

\subsection{Simulation Results}
The simulation waveform confirms correct behavior:
\begin{itemize}
    \item With switches = 3: LEDs display \texttt{0003 $\rightarrow$ 0002 $\rightarrow$ 0001 $\rightarrow$ 0000}
    \item With switches = 5: LEDs display \texttt{0005 $\rightarrow$ 0004 $\rightarrow$ 0003 $\rightarrow$ 0002 $\rightarrow$ 0001 $\rightarrow$ 0000}
    \item The processor returns to the WAIT state after each countdown completes.
\end{itemize}

\textit{(Attach simulation waveform screenshot here)}

% ============================================================
\section{Part B: Instruction Extension}
% ============================================================

\subsection{Objective}
Implement three new RISC-V instructions of different types that are not already supported by the processor.

\subsection{Selected Instructions}

\begin{table}[h]
\centering
\caption{New Instructions Implemented}
\begin{tabular}{@{}llll@{}}
\toprule
\textbf{Instruction} & \textbf{Type} & \textbf{Opcode} & \textbf{Description} \\
\midrule
SLT  & R-type & \texttt{0110011} & Set Less Than (signed comparison) \\
BGE  & B-type & \texttt{1100011} & Branch if Greater or Equal (signed) \\
LUI  & U-type & \texttt{0110111} & Load Upper Immediate \\
\bottomrule
\end{tabular}
\end{table}

% ----- SLT -----
\subsection{Instruction 1: SLT (R-type)}

\subsubsection{Functionality}
\texttt{SLT rd, rs1, rs2} sets \texttt{rd = 1} if \texttt{rs1 < rs2} (signed comparison), otherwise \texttt{rd = 0}. The encoding uses funct3 = \texttt{010} and funct7 = \texttt{0000000}.

\subsubsection{Hardware Modifications}

\textbf{ALU.v} --- Added one new case for the SLT operation:
\begin{lstlisting}[style=verilog, caption={ALU modification for SLT}]
// Added inside the case(ALUControl) block:
4'b0111: ALUResult = ($signed(A) < $signed(B)) ? 32'd1 : 32'd0;
\end{lstlisting}

\textbf{ALUControl.v} --- Added routing for funct3 = \texttt{010} under ALUOp = \texttt{10}:
\begin{lstlisting}[style=verilog, caption={ALUControl modification for SLT}]
// Added inside the ALUOp 2'b10 case:
3'b010 : ALUControl = 4'b0111; // SLT
\end{lstlisting}

\subsubsection{Datapath Flow}
The instruction flows through the standard R-type path: the opcode \texttt{0110011} sets ALUOp = \texttt{10} and RegWrite = 1. The ALUControl decodes funct3 = \texttt{010} to generate ALUControl = \texttt{0111}, which triggers the signed comparison in the ALU. The result (0 or 1) is written back to the destination register.

% ----- BGE -----
\subsection{Instruction 2: BGE (B-type)}

\subsubsection{Functionality}
\texttt{BGE rs1, rs2, offset} branches to \texttt{PC + offset} if \texttt{rs1 $\geq$ rs2} (signed). The encoding uses funct3 = \texttt{101} with the branch opcode \texttt{1100011}.

\subsubsection{Hardware Modifications}

\textbf{TopLevelProcessor.v} --- Extended the branch logic to include BGE:
\begin{lstlisting}[style=verilog, caption={Branch logic modification for BGE}]
wire branchTaken = Branch & (
    (funct3 == 3'b000 & Zero)   |        // BEQ
    (funct3 == 3'b001 & ~Zero)  |        // BNE
    (funct3 == 3'b101 & ~ALUResult[31])  // BGE (NEW)
);
\end{lstlisting}

\subsubsection{Datapath Flow}
BGE shares the branch opcode, so MainControl sets Branch = 1 and ALUOp = \texttt{01}. The ALUControl outputs SUB (\texttt{0001}), so the ALU computes \texttt{rs1 - rs2}. If the result's sign bit (bit 31) is 0, the subtraction is non-negative, meaning \texttt{rs1 $\geq$ rs2}. The condition \texttt{\~{}ALUResult[31]} captures this, and the branch is taken.

% ----- LUI -----
\subsection{Instruction 3: LUI (U-type)}

\subsubsection{Functionality}
\texttt{LUI rd, imm} loads a 20-bit immediate into the upper 20 bits of \texttt{rd}, with the lower 12 bits set to zero. The result is \texttt{rd = imm $\ll$ 12}.

\subsubsection{Hardware Modifications}

\textbf{MainControl.v} --- Added \texttt{LUI} output signal and a new case:
\begin{lstlisting}[style=verilog, caption={MainControl modification for LUI}]
// New output signal added:
output reg LUI

// New case in opcode decoder:
7'b0110111: begin // LUI
    RegWrite = 1'b1;
    LUI = 1'b1;
end
\end{lstlisting}

\textbf{TopLevelProcessor.v} --- Added LUI to the write-back multiplexer:
\begin{lstlisting}[style=verilog, caption={Write-back modification for LUI}]
assign regWriteData = (Jump | JumpReg) ? PCplus4
                    : (LUI ? imm : writeBackData);
\end{lstlisting}

\subsubsection{Datapath Flow}
LUI uses a unique opcode \texttt{0110111}. The MainControl sets RegWrite = 1 and LUI = 1. The immediate generator (immGen) already supports this opcode and produces \texttt{\{instruction[31:12], 12'b0\}}. The LUI signal routes this immediate value directly to the register write-back path, bypassing both the ALU and the memory. The immediate is written directly into the destination register.

\subsection{Verification Test Program}
A test program was written to verify all three instructions:

\begin{lstlisting}[style=asm, caption={Test program for new instructions}]
addi x6, x0, 512      # LED address

# TEST 1: SLT
addi x7, x0, 5        # x7 = 5
addi x8, x0, 10       # x8 = 10
slt  x9, x7, x8       # x9 = (5 < 10) = 1
sw   x9, 0(x6)         # LEDs show 0x0001

slt  x9, x8, x7       # x9 = (10 < 5) = 0
sw   x9, 0(x6)         # LEDs show 0x0000

# TEST 2: BGE
addi x7, x0, 10       # x7 = 10
addi x8, x0, 5        # x8 = 5
addi x9, x0, 0xAA     # x9 = 170
bge  x7, x8, skip     # 10 >= 5, branch (skip next)
addi x9, x0, 0xFF     # SKIPPED if BGE works
skip: sw x9, 0(x6)    # LEDs show 0x00AA (not 0xFF)

# TEST 3: LUI
lui  x9, 0xABCDE      # x9 = 0xABCDE000
sw   x9, 0(x6)         # LEDs show 0xE000 (lower 16 bits)

halt: jal x0, halt     # infinite loop
\end{lstlisting}

\subsection{Simulation Results}
The simulation confirms correct operation of all three instructions:
\begin{table}[h]
\centering
\caption{Part B Test Results}
\begin{tabular}{@{}llll@{}}
\toprule
\textbf{Test} & \textbf{Operation} & \textbf{Expected LEDs} & \textbf{Observed LEDs} \\
\midrule
SLT (true)  & \texttt{slt x9, 5, 10}    & \texttt{0x0001} & \texttt{0x0001} \\
SLT (false) & \texttt{slt x9, 10, 5}    & \texttt{0x0000} & \texttt{0x0000} \\
BGE         & \texttt{bge 10, 5, skip}   & \texttt{0x00AA} & \texttt{0x00AA} \\
LUI         & \texttt{lui x9, 0xABCDE}   & \texttt{0xE000} & \texttt{0xE000} \\
\bottomrule
\end{tabular}
\end{table}

\textit{(Attach simulation waveform screenshot here)}

% ============================================================
\section{Part C: Program Demonstration -- Fibonacci Series}
% ============================================================

\subsection{Objective}
Demonstrate a meaningful assembly program executing on the processor with correct stack usage for procedure calls.

\subsection{Program Description}
The program computes the first 10 Fibonacci numbers (1, 1, 2, 3, 5, 8, 13, 21, 34, 55) and displays each on the LEDs. It consists of two parts:

\begin{itemize}
    \item \textbf{Main loop (PC 0--52):} Initializes registers (\texttt{prev=0}, \texttt{curr=1}, \texttt{count=10}), then iterates: calls the display subroutine, computes the next Fibonacci number (\texttt{next = prev + curr}), shifts values, decrements the counter, and loops.
    \item \textbf{Display subroutine (PC 72--112):} A proper function that uses the stack to save and restore \texttt{ra} (return address) and \texttt{a0} (argument). It writes the Fibonacci number to the LEDs, runs a delay loop, then restores registers and returns.
\end{itemize}

\subsection{Register Usage}
\begin{table}[h]
\centering
\caption{Register Allocation for Fibonacci Program}
\begin{tabular}{@{}lll@{}}
\toprule
\textbf{Register} & \textbf{Purpose} & \textbf{Convention} \\
\midrule
\texttt{x2 (sp)}  & Stack pointer (initialized to 511) & Callee-saved \\
\texttt{x6 (t1)}  & LED address (512)                  & Temporary \\
\texttt{x7 (t2)}  & Previous Fibonacci number          & Temporary \\
\texttt{x8 (s0)}  & Current Fibonacci number           & Temporary \\
\texttt{x9 (s1)}  & Next Fibonacci number (temp)       & Temporary \\
\texttt{x10 (a0)} & Argument to subroutine             & Argument \\
\texttt{x13 (a3)} & Delay counter                      & Temporary \\
\texttt{x14 (a4)} & Loop counter (10 iterations)       & Temporary \\
\texttt{x1 (ra)}  & Return address                     & Saved on stack \\
\bottomrule
\end{tabular}
\end{table}

\subsection{Assembly Listing}
\begin{lstlisting}[style=asm, caption={Fibonacci Program}]
# === INITIALIZATION ===
addi x2, x0, 511      # sp = 511 (top of data memory)
addi x6, x0, 512      # x6 = LED address
addi x7, x0, 0        # prev = 0
addi x8, x0, 1        # curr = 1
addi x14, x0, 10      # count = 10

# === MAIN LOOP ===
loop:
    add  x10, x8, x0   # a0 = curr (argument for subroutine)
    jal  x1, display    # call display_and_delay
    add  x9, x7, x8    # next = prev + curr
    add  x7, x8, x0    # prev = curr
    add  x8, x9, x0    # curr = next
    addi x14, x14, -1  # count--
    bne  x14, x0, loop # if count != 0, loop

# === DONE ===
    sw   x0, 0(x6)     # clear LEDs
    jal  x0, halt       # halt (infinite loop)

# === DISPLAY_AND_DELAY SUBROUTINE ===
display:
    addi x2, x2, -8    # push stack (room for 2 words)
    sw   x1, 4(x2)     # save return address
    sw   x10, 0(x2)    # save a0
    sw   x10, 0(x6)    # write a0 to LEDs
    addi x13, x0, 5    # delay counter = 5
delay:
    addi x13, x13, -1  # delay--
    bne  x13, x0, delay # loop until zero
    lw   x10, 0(x2)    # restore a0
    lw   x1, 4(x2)     # restore return address
    addi x2, x2, 8     # pop stack
    jalr x0, 0(x1)     # return to caller
\end{lstlisting}

\subsection{Stack Usage}
The subroutine demonstrates correct RISC-V calling convention:
\begin{enumerate}
    \item \textbf{Prologue:} \texttt{sp} is decremented by 8 bytes to make room for two 32-bit words. The return address (\texttt{ra}) is saved at \texttt{sp+4} and the argument (\texttt{a0}) at \texttt{sp+0}.
    \item \textbf{Body:} The subroutine writes to LEDs and runs a delay loop.
    \item \textbf{Epilogue:} \texttt{a0} and \texttt{ra} are restored from the stack. \texttt{sp} is incremented by 8 to pop the frame. The subroutine returns via \texttt{jalr}.
\end{enumerate}

\subsection{Simulation Results}
The simulation waveform confirms the Fibonacci sequence on the LEDs:

\begin{table}[h]
\centering
\caption{Fibonacci Execution Results}
\begin{tabular}{@{}cl@{}}
\toprule
\textbf{Iteration} & \textbf{LED Value} \\
\midrule
1  & 1  \\
2  & 1  \\
3  & 2  \\
4  & 3  \\
5  & 5  \\
6  & 8  \\
7  & 13 \\
8  & 21 \\
9  & 34 \\
10 & 55 \\
End & 0 (cleared) \\
\bottomrule
\end{tabular}
\end{table}

\textit{(Attach simulation waveform screenshot here)}

% ============================================================
\section{Summary of Modified Files}
% ============================================================

\begin{table}[h]
\centering
\caption{Files Modified and Reason}
\begin{tabular}{@{}lll@{}}
\toprule
\textbf{File} & \textbf{Modification} & \textbf{For Instruction} \\
\midrule
ALU.v              & Added case \texttt{4'b0111} for signed comparison & SLT \\
ALUControl.v       & Added funct3 \texttt{3'b010} $\rightarrow$ \texttt{4'b0111} mapping & SLT \\
ALUControl.v       & Fixed encoding to match ALU operations & All (bug fix) \\
MainControl.v      & Added \texttt{LUI} output signal and opcode case & LUI \\
TopLevelProcessor.v & Added BGE condition to \texttt{branchTaken} & BGE \\
TopLevelProcessor.v & Added LUI write-back path (imm $\rightarrow$ rd) & LUI \\
instructionMemory.v & Updated with test/Fibonacci programs & Part B, C \\
\bottomrule
\end{tabular}
\end{table}

% ============================================================
\section{Conclusion}
% ============================================================

This project successfully demonstrated the integration of RISC-V assembly programs with a hardware-implemented single-cycle processor. The Lab 10 FSM program was executed on the processor from Lab 11, verifying correct instruction sequencing and memory-mapped I/O operation. Three new instructions (SLT, BGE, LUI) of different types were implemented with minimal hardware modifications and verified through simulation. A Fibonacci series program with proper stack-based subroutine calls was designed and executed, producing the correct output sequence on the LEDs. All programs were verified through behavioral simulation and are ready for FPGA demonstration.

\end{document}
